\documentclass[11pt]{article}
\usepackage{amsmath,amssymb,amsthm,mathtools}
\usepackage[margin=1in]{geometry}\usepackage{url}
\newtheorem{theorem}{Theorem}\newtheorem{lemma}[theorem]{Lemma}\newtheorem{proposition}[theorem]{Proposition}
\theoremstyle{remark}\newtheorem{remark}[theorem]{Remark}
\newcommand{\pb}{\overline{p}}
\newcommand{\leg}[2]{\left(\frac{#1}{#2}\right)}
\title{Overpartition congruences modulo 17, 19 and 23:\\ proofs of three conjectures of Wei}
\author{Anonymous}\date{12 October 2026}
\setlength{\emergencystretch}{3em}
\begin{document}\maketitle
\begin{abstract}
Wei recently proved Ramanujan-type congruences for the overpartition function $\pb(n)$ modulo $11$ and $13$, and stated three further congruences. Modulo $17$ and $23$ they were expected to follow from his method after Sturm-bound computations involving more than $10^9$ coefficients, which were not completed. Modulo $19$ the congruence was stated to be beyond the methods of the paper. We prove all three. For the moduli $17$ and $23$ we compute the Hecke operator $T(\ell^2)$ on a space of dimension at most $6$ and isolate a correction term that vanishes on the relevant residue class; the computation needs fewer than $4000$ coefficients. For the modulus $19$ the relevant generating function is only weakly holomorphic. We multiply it by two eta-quotients congruent to $1$ modulo $19$ before projecting onto the required classes, and show that the projection is a cusp form of weight $35/2$ on $\Gamma_0(23104)$. The congruence then follows from a $T(17^2)$ eigen-relation checked up to the Sturm bound.
\end{abstract}

\section{Introduction}
Let $\pb(n)$ denote the number of overpartitions of $n$~\cite{CL}, so that $\sum_{n\ge0}\pb(n)q^n=1/\varphi(-q)$ with $\varphi(q)=\sum_{k\in\mathbb Z}q^{k^2}$. Wei~\cite{Wei} proved infinite families of congruences for $\pb$ modulo $11$ and $13$ by Hecke operators on the spaces $M_k(\Gamma_0(4))$, which are spanned by products of $\varphi$ and $F(q)=\sum_{n\ \mathrm{odd}}\sigma(n)q^n$. In his closing section he recorded the following.

\begin{theorem}\label{thm:1723}
For all $n\ge0$,
\begin{align*}
\pb\bigl(17\cdot11^2(88n+m)\bigr)&\equiv0\pmod{17}, && m\in\{19,35,43,51,83\},\\
\pb\bigl(23\cdot13^2(104n+l)\bigr)&\equiv0\pmod{23}, && l\in\{29,53,61,69,77,101\}.
\end{align*}
\end{theorem}

\begin{theorem}\label{thm:19}
$\pb(17^2\cdot4n)\equiv0\pmod{19}$ for every integer $n\ge0$ with $n\equiv3\pmod 8$, $\leg n{17}=-1$ and $\leg n{19}=1$.
\end{theorem}

Wei states the congruences of Theorem~\ref{thm:1723}, explains that they reduce to Sturm-bound checks of about $1.2\cdot10^9$ and $4.0\cdot10^9$ coefficients on $\Gamma_1(340736)$ and $\Gamma_1(562432)$, and notes that these computations had not been completed. He states Theorem~\ref{thm:19} as a conjecture that cannot be proved by the methods of~\cite{Wei}. A further conjecture of~\cite{Wei}, modulo $7$, is treated in the modulo-$7$ paper in this versioned collection.

\paragraph{Contribution.} For Theorem~\ref{thm:1723}, Wei had already reduced the statements to finite computations and expected them to follow from his method. Our contribution there is a short structural argument replacing a computation of more than $10^9$ coefficients by one with fewer than $4000$. Theorem~\ref{thm:19} settles a conjecture that~\cite{Wei} places beyond its methods.

\paragraph{Relation to general results.} Treneer's pole-removal method~\cite{Treneer} is a standard input. Ryan, Sirolli, Villegas-Morales and Zheng~\cite[Theorems 1.3 and 1.5]{RSVZ} already give square-prime, Legendre-restricted congruences as well as cube-prime families. A cube-versus-square distinction would therefore be incorrect. The comparison must use the actual parameters: our modulo-$19$ target arguments are prime to $19$, whereas their displayed modulo-$19$ families contain that factor; the modulo-$17/23$ targets have the specified small primes and additional classes modulo $8$. Wei already supplies their starting forms in his closing section. The remaining contribution is the stated proof and computational reduction, not the invention of those forms or of square-prime congruences. The consulted sources did not supply these exact conclusions; this is a bounded priority assessment.

The residue classes in Theorem~\ref{thm:1723} are exactly the $N$ with $N\equiv3\pmod 8$ and $\leg N{11}=-1$, respectively $N\equiv5\pmod8$ and $\leg N{13}=1$. This was checked by enumeration modulo $88$ and $104$.

\paragraph{Method.} Both proofs use Shimura's operators $T(\ell^2)$ on half-integral weight forms~\cite{Shimura,Ono}. For $f=\sum a(n)q^n\in M_{\lambda+1/2}(\Gamma_0(N))$ (trivial character, $4\mid N$, $\ell\nmid N$),
\begin{equation}\label{eq:T}
(f\mid T(\ell^2))(n)=a(\ell^2n)+\leg{(-1)^\lambda n}{\ell}\ell^{\lambda-1}a(n)+\ell^{2\lambda-1}a(n/\ell^2),
\end{equation}
with $a(x)=0$ for $x\notin\mathbb Z$. Theorem~\ref{thm:1723} needs only the low-dimensional spaces $M_{\lambda+1/2}(\Gamma_0(4))$. Theorem~\ref{thm:19} concerns $\pb(4n)$, whose generating function $\varphi(q)^3/\varphi(-q)^4$ has a pole at the cusp $0$. The new ingredient is a holomorphic representative modulo~$19$ of the relevant projection.

\paragraph{Sturm's theorem.} If $g\in M_k(\Gamma_0(N),\chi)$ has $p$-integral coefficients and $a(n)\equiv0\pmod p$ for $n\le\frac{k}{12}[\mathrm{SL}_2(\mathbb Z):\Gamma_0(N)]$, then $g\equiv0\pmod p$~\cite{Sturm}. For half-integral weight we apply it to $g^2$ (or $g\varphi$), which has integral weight. If the first nonzero coefficient of $g$ modulo $p$ is at $n_0$, then that of $g^2$ is at $2n_0$, because $\mathbb F_p[[q]]$ is an integral domain.

\section{Proof of Theorem~\ref{thm:1723}}
Write $(p,\ell,\lambda,r)=(17,11,7,3)$ or $(23,13,10,5)$.

\begin{lemma}[Wei's reduction]\label{lem:c}
There is $c=\sum c(n)q^n\in M_{\lambda+1/2}(\Gamma_0(4))$ with integer coefficients and $\sum\pb(pn)(-q)^n\equiv c\pmod p$. Explicitly, with $F$ as above,
\begin{align*}
c&=13F^2\varphi^7+13F\varphi^{11}+\varphi^{15} && (p=17),\\
c&=20F^5\varphi+17F^4\varphi^5+14F^3\varphi^9+5F^2\varphi^{13}+9F\varphi^{17}+\varphi^{21} && (p=23).
\end{align*}
\end{lemma}
\begin{proof}
We have $\sum\pb(n)(-q)^n=1/\varphi(q)$ and $\varphi(q)^p\equiv\varphi(q^p)\pmod p$, so $\varphi(q^p)\sum\pb(n)(-q)^n\equiv\varphi^{p-1}$. Extracting the exponents divisible by $p$ and replacing $q^p$ by $q$ gives $\varphi(q)\sum\pb(pn)(-q)^n\equiv\varphi^{p-1}\mid T(p)\pmod p$, where $T(p)$ is the Hecke operator on $M_{(p-1)/2}(\Gamma_0(4),\chi_{-4}^{(p-1)/2})$ (its second term carries $p^{(p-3)/2}\equiv0$). We computed $\varphi^{p-1}\mid T(p)$ exactly in the triangular basis $F^a\varphi^{p-1-4a}$. Its coordinates modulo $p$ are $(1,13,13,0,0)$ for $p=17$ and $(1,9,5,14,17,20)$ for $p=23$. Each is $\varphi$ times the stated $c$. For $p=17$ the coordinate of $F^4$ is $0$, so no $\varphi^{-1}$ term arises. Since $\varphi\in1+2q\mathbb Z[[q]]$ is a unit modulo $p$, the claim follows.
\end{proof}

\begin{proposition}\label{prop:h}
Modulo $p$, $c\mid T(\ell^2)\equiv\ell^{\lambda-1}c+h$ with $\ell^{\lambda-1}\equiv8\pmod{17}$, respectively $3\pmod{23}$, and
\begin{align*}
h&=12(F\varphi^{11}+F^2\varphi^7) && (p=17),\\
h&=16F\varphi^{17}+16F^2\varphi^{13}+14F^3\varphi^9+7F^4\varphi^5+15F^5\varphi && (p=23).
\end{align*}
\end{proposition}
\begin{proof}
$T(\ell^2)$ preserves $M_{\lambda+1/2}(\Gamma_0(4))$, which has the upper-triangular basis $F^a\varphi^{2\lambda+1-4a}$ of dimension $4$, respectively $6$. We computed the integer matrix of $T(\ell^2)$ in this basis from~\eqref{eq:T}, checking membership of each image in the span on all available coefficients. We then reduced $c\mid T(\ell^2)-\ell^{\lambda-1}c$ modulo $p$.
\end{proof}

\begin{lemma}\label{lem:hclass}
$h(n)\equiv0\pmod p$ for all $n\equiv r\pmod 8$.
\end{lemma}
\begin{proof}
Let $g=\sum_{n\equiv r\ (8)}h(n)q^n$. Since $r$ is odd, averaging the twists of $h$ by the four real characters modulo $8$ shows $g\in M_{\lambda+1/2}(\Gamma_0(256))$~\cite[Prop.~2.4 and Remark~2.5]{Wei}. Then $g\varphi\in M_{\lambda+1}(\Gamma_0(256),\chi_{-4}^{\lambda+1})$ has integer coefficients. The Sturm bound is $(\lambda+1)\cdot384/12=256$, respectively $352$. The coefficients of $g\varphi$ vanish modulo $p$ for $n<4000$, so $g\varphi\equiv0$, and $g\equiv0$ because $\varphi$ is a unit.
\end{proof}

\begin{proof}[Proof of Theorem~\ref{thm:1723}]
Let $N\equiv r\pmod8$ with $\leg N\ell=-1$ (for $p=17$), respectively $+1$ (for $p=23$). Then $\ell\nmid N$, and the character value in~\eqref{eq:T} is $\leg{-1}{11}\leg N{11}=(-1)(-1)=1$, respectively $\leg N{13}=1$ since $\lambda=10$ is even. Comparing $N$-th coefficients in Proposition~\ref{prop:h},
\[c(\ell^2N)+\ell^{\lambda-1}c(N)\equiv\ell^{\lambda-1}c(N)+h(N)\pmod p,\]
so $c(\ell^2N)\equiv h(N)\equiv0$ by Lemma~\ref{lem:hclass}. By Lemma~\ref{lem:c}, $\pb(p\ell^2N)\equiv(-1)^{\ell^2N}c(\ell^2N)\equiv0\pmod p$.
\end{proof}

\section{Proof of Theorem~\ref{thm:19}}
Let $c(n)=\pb(4n)$. A $2$-dissection gives
\[F_4:=\sum_{n\ge0}c(n)q^n=\frac{\varphi(q)^3}{\varphi(-q)^4}=\frac{\eta(2z)^{19}}{\eta(z)^{14}\eta(4z)^{6}},\]
The identity follows algebraically, not from a finite check: with $E(\sum a_nq^n)=\sum a_{2n}q^n$, the classical relations $\varphi(q)\varphi(-q)=\varphi(-q^2)^2$, $E\varphi=\varphi(q^2)$ and $E(\varphi^2)=\varphi^2$ give successively
\[E(1/\varphi(-q))=\varphi(q^2)/\varphi(-q)^2,
\qquad E(\varphi(q^2)\varphi(q)^2/\varphi(-q^2)^4)=F_4.\]
Since $\varphi(q)$ and $\varphi(-q)=2\varphi(q^4)-\varphi(q)$ transform under $\Gamma_0(16)$ with Shimura's theta multiplier $j(\gamma,z)$, $F_4$ is weakly holomorphic of weight $-1/2$. Its raw exponent at odd-denominator cusps is $-1/4$ (multiply by cusp width for the local-parameter order). Theorem~\ref{thm:19} asserts $c(289n)\equiv0\pmod{19}$ on the stated classes.

\paragraph{A representative without poles at most cusps.} Put $E_1=\eta(z)^{19}/\eta(19z)$ and $E_2=\eta(19z)^{19}/\eta(361z)$. Both lie in $1+19q\mathbb Z[[q]]$, because $(1-q^n)^{19}\equiv1-q^{19n}$. By the criterion of Gordon--Hughes--Newman and Ligozat, $E_1E_2\in M_{18}(\Gamma_0(361))$ with trivial character. Set
\[F''=F_4E_1E_2=\frac{\eta(z)^{5}\eta(2z)^{19}\eta(19z)^{18}}{\eta(4z)^{6}\eta(361z)}\equiv F_4\ (\mathrm{mod}\ 19).\]
Then $F''(\gamma z)=j(\gamma,z)^{35}F''(z)$ on $\Gamma_0(16\cdot361)$. Indeed $F_4(\gamma z)=j(\gamma,z)^{-1}F_4(z)$, and $(E_1E_2)(\gamma z)=(cz+d)^{18}(E_1E_2)(z)$. Since $j(\gamma,z)^2=\leg{-1}d(cz+d)$, we have $(cz+d)^{18}=j(\gamma,z)^{36}$. So $F''$ is weakly holomorphic of weight $35/2$ with trivial character. Multiplying by an eta-quotient congruent to $1$ modulo $p$ to remove poles is the standard device of Treneer~\cite{Treneer}. The point here is that the multiplication must precede the projection. At a cusp $a/c$ with $361\nmid c$, the factor $E_1$ (if $19\nmid c$) or $E_2$ (if $19\parallel c$) has order large enough to cancel the pole of $F_4$. So $F''$ has positive order at every cusp with $361\nmid c$.

\paragraph{Projection.} Let $G''=\sum c''(n)q^n$, summed over $n\equiv3\pmod8$ with $\leg n{19}=1$, where $c''(n)$ are the coefficients of $F''$. Then $G''\in\mathbb Z[[q]]$ and $G''\equiv\sum_{n\equiv3\,(8),\,(n/19)=1}c(n)q^n\pmod{19}$. Explicitly $G''(z)=\sum_{t\bmod152}W_tF''(z+t/152)$. Writing $t=19j+8k$, the weights are $W_t=\frac18e(-3j/8)\,b_k$ with $b_0=9/19$ and $b_k=-(1+\chi(k)g)/38$, where $\chi=\leg\cdot{19}$ and $g$ is the Gauss sum. These weights reproduce the indicator function exactly.

\begin{lemma}
$G''(\gamma z)=j(\gamma,z)^{35}G''(z)$ for all $\gamma\in\Gamma_0(23104)$, where $23104=64\cdot361$.
\end{lemma}
\begin{proof}
Let $\gamma=\begin{pmatrix}a&b\\c&d\end{pmatrix}\in\Gamma_0(23104)$, so $152^2\mid c$. For each $t$ put $t'\equiv td^2\pmod{152}$. Then $\gamma'=T_{t/152}\gamma T_{-t'/152}$ is in $\Gamma_0(16\cdot361)$ with lower-right entry $d'=d-ct'/152\equiv d$ modulo $8$ and modulo $c/23104$. Hence $j(\gamma',z+t'/152)=j(\gamma,z)$. Since $d^2\equiv1\pmod 8$ and $\chi(kd^2)=\chi(k)$, we have $W_{t'}=W_t$, and the sum is permuted.
\end{proof}

\begin{proposition}\label{prop:cusp}
$G''$ is a cusp form: $G''\in S_{35/2}(\Gamma_0(23104))$.
\end{proposition}
\begin{proof}
At the $228$ cusps with $361\nmid c$, every translate is expanded at a cusp whose reduced denominator is again not divisible by $361$; exact order arithmetic gives strictly positive order in all $228\cdot152=34656$ cases. At the other twelve cusps, some translated summands have nonpositive exponents; this does not say that $F''$ itself has a pole at each of them. Appendix~\ref{app:cusp} specifies the complete finite obligation and expansion convention. Grouping by exponent and rational square class leaves $28$ cyclotomic sums with $32928$ summands. Every sum reduces exactly to zero modulo its cyclotomic polynomial. Thus all principal parts and constant terms vanish. This exact calculation, not numerical decay, proves the assertion.

Historical numerical corroboration used a separately written program; it is not the exact proof or a newly repeated scan. It evaluated $|G''(\gamma z)(cz+d)^{-35/2}|$ at $200$-digit precision at all $240$ cusps, along vertical rays of increasing height. It observed exponential decay everywhere, with stable fitted orders of $16$, $19$, $4$ and $11$ in the local parameter, and found no constant term or pole. At the $12$ critical cusps the cancellation is by factors $10^{-63}$ to $10^{-118}$, well within the working precision. The same program checked the transformation law on several matrices of $\Gamma_0(23104)$.
\end{proof}

\begin{proof}[Proof of Theorem~\ref{thm:19}]
Since $17\nmid23104$, $T(289)$ acts on $S_{35/2}(\Gamma_0(23104))$ by~\eqref{eq:T} with $\lambda=17$. Moreover $17^{16}\equiv5$ and $17^{33}\equiv17^{15}\pmod{19}$. Let $H=G''\mid T(289)-14G''$, a cusp form with integer coefficients. Its square belongs to $S_{35}(\Gamma_0(23104),\chi_{-4})$. The index is $36480$, so the character-version Sturm bound is $106400$. It suffices to check $H(n)\equiv0\pmod{19}$ for $0\le n\le53200$, inclusively. One implementation generates $F_4$ through exponent $289\cdot53200=15374800$ (with ten spare coefficients); a separate implementation generates $\pb$ through argument $4\cdot289\cdot53200=61499200$. Both test the same finite obligation, with no exception. Hence, for all $n$ in the support,
\[c(289n)+5\leg n{17}c(n)+17^{15}c(n/289)\equiv14\,c(n)\pmod{19}.\]
If $n\equiv3\pmod8$, $\leg n{19}=1$ and $\leg n{17}=-1$, then $17\nmid n$, $289n$ lies in the same classes, and the relation gives $c(289n)\equiv19c(n)\equiv0$.
\end{proof}

\begin{remark}
The scalar $14$ differs from $1+17^{33}\equiv8\pmod{19}$; no particular Eisenstein series is being identified here. The relation does not extend to arbitrary coefficients outside the projection. Projecting $F_4$ before multiplying by $E_1E_2$ leaves poles, so the order of the two operations matters.
\end{remark}

\section{Computations}
All computations use exact integer or modular arithmetic (python-flint), except the high-precision confirmation in Proposition~\ref{prop:cusp}. The scripts and logs are in the accompanying evidence bundle. Each theorem was also confirmed directly: Theorem~\ref{thm:1723} for arguments up to $2\cdot10^7$ ($552$ and $295$ cases), and Theorem~\ref{thm:19} for all $n\le53200$.

\appendix
\section{Exact cusp certificate}\label{app:cusp}
Write $e(x)=\exp(2\pi ix)$. The eta-exponent vector, in increasing argument order, is
\[(\delta,r_\delta)=(1,5),(2,19),(4,-6),(19,18),(361,-1).\]
Represent the cusps by $a/c$ with $c\mid L=23104$, $\gcd(a,c)=1$, and one numerator in each unit class modulo $\gcd(c,L/c)$. There are $240$ such cusps. Choose $A_s=\left(\begin{smallmatrix}a&b\\c&d\end{smallmatrix}\right)$ with $ad-bc=1$. The expansion convention is
\[(-i(c\tau+d))^{-35/2}G''(A_s\tau)=\sum_E C_E e(E\tau).
\]
The cusp width is $w=L/\gcd(c^2,L)$, so the exponent in the width-normalised local parameter is $wE$. Multiplication by a nonzero phase to change slash conventions does not change vanishing.

\begin{center}\small
\begin{tabular}{r r r l r}
$c$&cusps&$w$&raw $E\le0$&groups\\\hline
361&1&64&$-1/4,0$&3\\
722&1&16&$-1,0$&3\\
1444&2&4&$-4,0$&6\\
2888&4&1&$-16,0$&12\\
5776&2&1&$0$&2\\
11552&1&1&$0$&1\\
23104&1&1&$0$&1\\\hline
Total&12&&&28
\end{tabular}
\end{center}
The listed exponents are those arising from translated summands before cancellation. All resulting group sums are zero.

\paragraph{Matrix and multiplier convention.}
For each $t=19j+8k$, decompose the integer matrix
\[M_t=\begin{pmatrix}152a+tc&152b+td\\152c&152d\end{pmatrix}
=\gamma\begin{pmatrix}u_1&u_2\\0&u_3\end{pmatrix},\qquad \gamma\in\mathrm{SL}_2(\mathbb Z),\ u_1,u_3>0.
\]
This is obtained by the extended Euclidean algorithm on the first column, with $u_1$ its positive gcd. For each eta factor similarly decompose $\operatorname{diag}(\delta,1)\gamma=\gamma_\delta V_\delta$. For a matrix $(a',b';c',d')$ with $c'>0$, use
\[\eta(\gamma_\delta z)=e\!\left(\frac{a'+d'}{24c'}-\frac{s(d',c')}{2}\right)
[-i(c'z+d')]^{1/2}\eta(z),\]
where $s$ is the rational Dedekind sum and the analytic square root is fixed on the upper half-plane. This is the convention in Kong--Teo~\cite[Eq.~(1.2)]{KT} after absorbing the $e(-1/8)$ phase into the square root.

\paragraph{Finite truncation and exact arithmetic.}
The decompositions give a leading exponent $E_0$, a root-of-unity phase, a positive square root of a rational number, and factors
\[\prod_{\delta}\prod_{n\ge1}(1-e(n\beta_\delta)q^{n\alpha_\delta})^{r_\delta},
\qquad \alpha_\delta>0.\]
If $E_0>0$, no term is needed. Otherwise expand only increments at most $-E_0$. All increments are positive; consequently omitted factors and higher binomial terms cannot contribute to an exponent $E\le0$. For $r<0$ use $\binom r m=(-1)^m\binom{-r+m-1}m$. All phases and amplitudes before square roots are rational and retained exactly.

Multiply by the specified projection weights $W_t$ and group by $E$ and rational square class. Within each group, choose the least common denominator $n$ of phases and form $P(X)=\sum_j q_j X^{nr_j\bmod n}\in\mathbb Q[X]$. Check $P\bmod\Phi_n=0$. The actual fields have $n=38,76,152$. This separate vanishing of all square-class groups is sufficient; no assumption of their linear independence is required. The certificate contains $32928$ summands in $28$ groups. Projection denominators divisible by $19$ cause no reduction problem: $G''$ is defined as retention of integer coefficients, and the Fourier weights are used only to prove its modularity and cusp behaviour in characteristic zero.

\paragraph{Reproduction.}
From the collection root run \texttt{python verify.py mod17-19-23}. The wrapper pins the lifted eta vector, checks all $34656$ positive orders, requires exactly twelve critical cusps and all $28$ zero groups, and performs the inclusive Hecke--Sturm check. It fails on an unmet obligation, not merely prints a Boolean. Numerical cusp decay is retained as corroboration only.

\begin{proposition}[Sufficient lift--project criterion]
Let $f$ be weakly holomorphic with integral Fourier coefficients, and let $E\equiv1\pmod p$ be an integral eta product or quotient. Suppose a finite Fourier projection $P$ retains specified integer coefficients, $P(Ef)$ has a proved modular transformation law at a finite-index level, and every nonpositive cusp coefficient vanishes by complete finite expansions as above. Then $P(Ef)$ is an integral cusp form congruent to $P(f)$ modulo $p$. Any proposed Hecke congruence for it follows from a valid Sturm check in the resulting weight, level and character.
\end{proposition}
\begin{proof} The transformation law and cusp expansions establish cusp membership. Coefficient retention commutes with reduction, so $P(Ef)\equiv P(f)$. Apply the relevant Hecke operator and Sturm theorem to their difference.\end{proof}
This is a sufficient certificate pattern, not an automatic existence theorem for a suitable $E$, nor a claim that this general principle is new.

\section*{Collection and verification}
This paper is part of the revised collection, version 1.1.0-candidate, DOI \url{https://doi.org/10.5281/zenodo.23319157}. Companion papers and exact verification code are archived together. The statements are unrefereed computer-assisted proof candidates; no formal verification or unaffiliated reproduction is asserted.

\begin{thebibliography}{99}
\bibitem{KT} Z.-Y.~Kong and L.-P.~Teo, An elementary proof of the transformation formula for the Dedekind eta function, arXiv:2302.03280 (2023), \url{https://arxiv.org/abs/2302.03280}.
\bibitem{CL} S.~Corteel and J.~Lovejoy, Overpartitions, \emph{Trans. Amer. Math. Soc.} 356 (2004), 1623--1635.
\bibitem{Ono} K.~Ono, \emph{The Web of Modularity: Arithmetic of the Coefficients of Modular Forms and $q$-series}, CBMS Regional Conf. Ser. Math. 102, Amer. Math. Soc., 2004.
\bibitem{Shimura} G.~Shimura, On modular forms of half integral weight, \emph{Ann. of Math.} 97 (1973), 440--481.
\bibitem{Sturm} J.~Sturm, On the congruence of modular forms, in \emph{Number Theory (New York, 1984--1985)}, Lecture Notes in Math. 1240, Springer, 1987, 275--280.
\bibitem{RSVZ} N.~C.~Ryan, N.~Sirolli, J.~C.~Villegas-Morales and Q.-Y.~Zheng, Explicit families of congruences for the overpartition function, \emph{Ramanujan J.} (2024); arXiv:2309.01792.
\bibitem{Treneer} S.~Treneer, Congruences for the coefficients of weakly holomorphic modular forms, \emph{Proc. London Math. Soc. (3)} 93 (2006), 304--324.
\bibitem{Wei} X.~Wei, New Ramanujan-type congruences for overpartitions modulo 11 and 13, arXiv:2603.08510 (2026).
\end{thebibliography}
\end{document}
